\section{Relation to the original manuscript and its formalisation}\label{app:notes}

The construction and almost all estimates are those of~\cite{Orig}. The
differences are as follows.

\begin{enumerate}[label=\arabic*.]
\item \emph{Many-row seed for all small errors.} In~\cite{Orig} the many-row
theorem assumes $\eta\le c/d^2$, because the cubic remainder of the
normalisation identity is bounded termwise by $O(t^3d)$, and the operator norm of
the noise by $t^2\fro{Z}^2\le1$. \Cref{lem:remainder} uses the first-order
constraint $X^TZ+Z^TX=0$ a second time to centre the weights, giving $O(t^3)$,
and the net bound controls $\opn Z$. \Cref{thm:manyrow} therefore holds for
$\eta\le c_*$.
\item \emph{No block partition.} With 1., the global argument needs no random
partition into blocks of size $\lceil Ld^6\rceil$ (and no uniform-subset variance
identity or averaging over block errors): the three regimes $d<D$, $n\ge Bd^2$,
$n<Bd^2$ of \cref{prop:equalrow} cover everything, since $A\log(2Bd^2)\le d$ for
$d\ge D$.
\item \emph{One toolbox.} The distance estimates, the potential, the static
balancing theorem, the local correction lemma and the dense-core Poisson
estimate are stated once (\cref{sec:toolbox}), and both seeds are routed
through \cref{def:seed} and \cref{cor:seed}.
\item \emph{Complete sampling estimates.} The probabilistic statements used by
the moderate seed are collected in \cref{lem:sample} and proved in full,
including the variance bookkeeping behind the dense core.
\end{enumerate}

\paragraph{Formal status.} The Lean~4 development accompanying~\cite{Orig}
proves \cref{thm:main,thm:projection} by the original route (with the block
partition and the $\eta\le c/d^2$ many-row theorem). The argument of
\cref{sec:manyrow,sec:assembly} here is not formalised; the other sections follow
the formalised proof up to constants and presentation.
